The prime numbers look random: no formula produces them, twin primes keep appearing, and so do long gaps. Yet in bulk
they obey precise laws. Their number up to $x$ is close to $x/\log x$; they are shared equally among the residue classes
that can contain them; the number of prime factors of a typical integer follows the bell curve of chance. This book proves
these laws, with all the earlier volumes working together. It builds Dirichlet's $L$-functions to find primes in progressions, proves the prime number theorem, studies
the zeros of the zeta function and Riemann's explicit formula, in which the primes are a superposition of waves, one
for each zero, and explains why primes of the form $4n+3$ win their race almost all the time. Then come sieves and
twin primes, the Erd\H{o}s--Kac theorem, random models of the primes, exponential sums, and a first taste of modular
forms. Every result is proved in full, each chapter opens with the story behind its ideas, and its Python notebook
turns the computer into a laboratory for primes.
